\documentclass[../../main.tex]{subfiles}
\begin{document}
\section{Lie 括号}\label{section:ch08-Lie括号}

本节引入一种将两个光滑向量场组合得到另一个向量场的重要方式。

设 $X$ 和 $Y$ 是光滑流形 $M$ 上的光滑向量场。给定光滑函数 $f\colon M\to\mathbb{R}$，可以将 $X$ 作用于 $f$ 得到另一个光滑函数 $Xf$（见\cref{proposition:ch08-14}）。反过来，可以将 $Y$ 作用于这个函数，得到又一个光滑函数 $YXf=Y(Xf)$。然而，运算 $f\mapsto YXf$ 一般不满足乘积法则，因而不能是向量场，如下例所示。

\begin{example}\label{example:ch08-24}
在 $\mathbb{R}^2$ 上定义向量场 $X=\partial/\partial x$ 和 $Y=x\partial/\partial y$，令 $f(x,y)=x$，$g(x,y)=y$。直接计算表明 $XY(fg)=2x$，而 $f(XY)g+g(XY)f=x$，故 $XY$ 不是 $C^\infty(\mathbb{R}^2)$ 的导子。
\end{example}

我们也可以以相反顺序应用相同的两个向量场，得到一个（通常不同的）函数 $XYf$。将这两个算子都应用于 $f$ 并相减，得到算子 $[X,Y]\colon C^\infty(M)\to C^\infty(M)$，称为 $X$ 和 $Y$ 的\textbf{Lie 括号}，定义为
\begin{align*}
[X,Y]f=XYf-YXf.
\end{align*}
关键事实是该算子是一个向量场。

\begin{lemma}\label{lemma:ch08-25}
任意一对光滑向量场的 Lie 括号是光滑向量场。
\end{lemma}

\begin{proof}
由\cref{proposition:ch08-15}，只需证明 $[X,Y]$ 是 $C^\infty(M)$ 的导子。对任意 $f,g\in C^\infty(M)$，计算
\begin{align*}
[X,Y](fg)&=X\bigl(Y(fg)\bigr)-Y\bigl(X(fg)\bigr)\\
&=X(fYg+gYf)-Y(fXg+gXf)\\
&=Xf\,Yg+f\,XYg+Xg\,Yf+g\,XYf\\
&\quad -Yf\,Xg-f\,YXg-Yg\,Xf-g\,YXf\\
&=f\,XYg+g\,XYf-f\,YXg-g\,YXf\\
&=f[X,Y]g+g[X,Y]f.
\end{align*}
\end{proof}

稍后在本章中我们将描述 Lie 括号的一个重要应用，在后续章节中将看到更多应用。遗憾的是，我们现在还无法给 Lie 括号一个几何解释，但将在第 9 章中这样做。现在先发展它们的一些基本性质。

向量场 $[X,Y]$ 在点 $p\in M$ 处的值是由公式
\begin{align*}
[X,Y]_pf=X_p(Yf)-Y_p(Xf)
\end{align*}
给出的 $p$ 处的导子。然而，该公式在计算中用处有限，因为它要求计算涉及 $f$ 的二阶导数的项，而这些项总是相互抵消。下一个命题给出 Lie 括号一个极其有用的坐标公式，其中的抵消已经被考虑在内。

\begin{proposition}\label{proposition:ch08-26}[Lie 括号的坐标公式]
设 $X,Y$ 是带边或不带边光滑流形 $M$ 上的光滑向量场，$X=X^i\partial/\partial x^i$ 和 $Y=Y^j\partial/\partial x^j$ 是 $X$ 和 $Y$ 在 $M$ 的某光滑局部坐标 $(x^i)$ 下的坐标表示。则 $[X,Y]$ 有如下坐标表示：
\begin{align}
[X,Y]=\left(X^i\frac{\partial Y^j}{\partial x^i}-Y^i\frac{\partial X^j}{\partial x^i}\right)\frac{\partial}{\partial x^j},\label{align:ch08-08}
\end{align}
或更简洁地，
\begin{align}
[X,Y]=(XY^j-YX^j)\frac{\partial}{\partial x^j}.\label{align:ch08-09}
\end{align}
\end{proposition}

\begin{proof}
因为我们已经知道 $[X,Y]$ 是光滑向量场，它对函数的作用由局部决定：$([X,Y]f)|_U=[X,Y](f|_U)$。因此只需在单个光滑坐标卡中计算，其中有
\begin{align*}
[X,Y]f&=X^i\frac{\partial}{\partial x^i}\left(Y^j\frac{\partial f}{\partial x^j}\right)-Y^j\frac{\partial}{\partial x^j}\left(X^i\frac{\partial f}{\partial x^i}\right)\\
&=X^i\frac{\partial Y^j}{\partial x^i}\frac{\partial f}{\partial x^j}+X^iY^j\frac{\partial^2 f}{\partial x^i\partial x^j}\\
&\quad -Y^j\frac{\partial X^i}{\partial x^j}\frac{\partial f}{\partial x^i}-Y^jX^i\frac{\partial^2 f}{\partial x^j\partial x^i}\\
&=X^i\frac{\partial Y^j}{\partial x^i}\frac{\partial f}{\partial x^j}-Y^j\frac{\partial X^i}{\partial x^j}\frac{\partial f}{\partial x^i},
\end{align*}
其中最后一步使用了光滑函数的混合偏导数可按任意顺序取这一事实。交换第二项中哑指标 $i$ 和 $j$ 的角色，得到\eqref{align:ch08-08}。
\end{proof}

\eqref{align:ch08-08}的一个平凡应用是计算任意光滑坐标卡中坐标向量场 $(\partial/\partial x^i)$ 的 Lie 括号：因为坐标向量场的分量函数全是常数，由此可知
\begin{align}
\left[\frac{\partial}{\partial x^i},\frac{\partial}{\partial x^j}\right]=0\quad\text{对所有 }i,j.\label{align:ch08-10}
\end{align}
（这也可由 Lie 括号的定义得到，本质上是光滑函数混合偏导数可交换这一事实的重述。）下面是一个稍微不那么平凡的计算。

\begin{example}\label{example:ch08-27}
由
\begin{align*}
X=x\frac{\partial}{\partial x}+\frac{\partial}{\partial y}+x(y+1)\frac{\partial}{\partial z},\qquad
Y=\frac{\partial}{\partial x}+y\frac{\partial}{\partial z}
\end{align*}
定义 $\mathfrak{X}(\mathbb{R}^3)$ 中的光滑向量场 $X,Y$。则\eqref{align:ch08-09}给出
\begin{align*}
[X,Y]&=X(1)\frac{\partial}{\partial x}+X(y)\frac{\partial}{\partial z}-Y(x)\frac{\partial}{\partial x}-Y(1)\frac{\partial}{\partial y}-Y\bigl(x(y+1)\bigr)\frac{\partial}{\partial z}\\
&=0\frac{\partial}{\partial x}+1\frac{\partial}{\partial z}-1\frac{\partial}{\partial x}-0\frac{\partial}{\partial y}-(y+1)\frac{\partial}{\partial z}\\
&=-\frac{\partial}{\partial x}-y\frac{\partial}{\partial z}.
\end{align*}
\end{example}

\begin{proposition}\label{proposition:ch08-28}[Lie 括号的性质]
Lie 括号对所有 $X,Y,Z\in\mathfrak{X}(M)$ 满足下列恒等式：
\begin{enumerate}[(a)]
\item \textbf{双线性}：对 $a,b\in\mathbb{R}$，
\begin{align*}
[aX+bY,Z]&=a[X,Z]+b[Y,Z],\\
[Z,aX+bY]&=a[Z,X]+b[Z,Y].
\end{align*}
\item \textbf{反称性}：
\begin{align*}
[X,Y]=-[Y,X].
\end{align*}
\item \textbf{Jacobi 恒等式}：
\begin{align*}
[X,[Y,Z]]+[Y,[Z,X]]+[Z,[X,Y]]=0.
\end{align*}
\item 对 $f,g\in C^\infty(M)$，
\begin{align}
[fX,gY]=fg[X,Y]+(fXg)Y-(gYf)X.\label{align:ch08-11}
\end{align}
\end{enumerate}
\end{proposition}

\begin{proof}
双线性和反称性是定义的明显推论。Jacobi 恒等式的证明只是一个计算：
\begin{align*}
&[X,[Y,Z]]f+[Y,[Z,X]]f+[Z,[X,Y]]f\\
&\quad=X[Y,Z]f-[Y,Z]Xf+Y[Z,X]f-[Z,X]Yf+Z[X,Y]f-[X,Y]Zf\\
&\quad=XYZf-XZYf-YZXf+ZYXf+YZXf-YXZf\\
&\qquad-ZXYf+XZYf+ZXYf-ZYXf-XYZf+YXZf.
\end{align*}
在最后这个表达式中所有项两两抵消。(d) 部分是由 Lie 括号定义的直接计算，留作练习。
\end{proof}

\begin{exercise}\label{exercise:ch08-29}
证明前述命题的 (d) 部分。
\end{exercise}

该命题 (d) 部分的意义此刻可能还不明显，但在下一章中将变得更清楚，在那里我们将看到它表达了 Lie 括号关于其两个参数都满足乘积法则这一事实（见\cref{corollary:ch09-39}）。

\begin{proposition}\label{proposition:ch08-30}[Lie 括号的自然性]
设 $F\colon M\to N$ 是带边或不带边流形之间的光滑映射，$X_1,X_2\in\mathfrak{X}(M)$，$Y_1,Y_2\in\mathfrak{X}(N)$ 是使得对 $i=1,2$ 有 $X_i$ 与 $Y_i$ 是 $F$-相关的向量场。则 $[X_1,X_2]$ 与 $[Y_1,Y_2]$ 是 $F$-相关的。
\end{proposition}

\begin{proof}
利用\cref{proposition:ch08-16}以及 $X_i$ 与 $Y_i$ 是 $F$-相关的这一事实，
\begin{align*}
X_1X_2(f\circ F)=X_1\bigl((Y_2f)\circ F\bigr)=(Y_1Y_2f)\circ F.
\end{align*}
类似地，
\begin{align*}
X_2X_1(f\circ F)=(Y_2Y_1f)\circ F.
\end{align*}
因此，
\begin{align*}
[X_1,X_2](f\circ F)&=X_1X_2(f\circ F)-X_2X_1(f\circ F)\\
&=(Y_1Y_2f)\circ F-(Y_2Y_1f)\circ F\\
&=\bigl([Y_1,Y_2]f\bigr)\circ F.
\end{align*}
\end{proof}

关于为什么这一性质称为"自然性"，见\cref{exercise:ch11-18}。当应用于特殊情形时，该结果有如下重要推论。首先考虑映射是微分同胚的情形。

\begin{corollary}\label{corollary:ch08-31}[Lie 括号的前推]
假设 $F\colon M\to N$ 是微分同胚，$X_1,X_2\in\mathfrak{X}(M)$。则 $F_*[X_1,X_2]=[F_*X_1,F_*X_2]$。
\end{corollary}

\begin{proof}
这正是\cref{proposition:ch08-30}在 $F$ 是微分同胚且 $Y_i=F_*X_i$ 时的特殊情形。
\end{proof}

第二个特殊情形是子流形的包含。

\begin{corollary}\label{corollary:ch08-32}[与子流形相切的向量场的括号]
设 $M$ 是光滑流形，$S$ 是 $M$ 中带边或不带边的浸入子流形。若 $Y_1$ 和 $Y_2$ 是 $M$ 上与 $S$ 相切的光滑向量场，则 $[Y_1,Y_2]$ 也与 $S$ 相切。
\end{corollary}

\begin{proof}
由\cref{proposition:ch08-23}，存在 $S$ 上的光滑向量场 $X_1$ 和 $X_2$，使得对 $i=1,2$ 有 $X_i$ 与 $Y_i$ 是 $\iota$-相关的（其中 $\iota\colon S\hookrightarrow M$ 是包含映射）。由\cref{proposition:ch08-30}，$[X_1,X_2]$ 与 $[Y_1,Y_2]$ 是 $\iota$-相关的，因此 $[Y_1,Y_2]$ 与 $S$ 相切。
\end{proof}

\end{document}